\documentclass[border=5mm]{standalone}
% ==============================================================================
% SETUP.TEX - CONFIGURACIÓN GLOBAL DEL PROYECTO
% ==============================================================================
% Este archivo centraliza todos los paquetes y estilos.
% Debe incluirse en el preámbulo del libro principal y de los archivos standalone.
% \PassOptionsToPackage{gray}{xcolor}
% --- 1. CODIFICACIÓN E IDIOMA ---
% fix-cm habilita las fuentes Computer Modern en tamaños arbitrarios (por debajo
% de 5 pt, entre otros). Debe cargarse antes que fontenc.
\usepackage{fix-cm}
\usepackage[utf8]{inputenc}
\usepackage[T1]{fontenc}
\usepackage[spanish, es-tabla]{babel}  
\usepackage{csquotes} % Recomendado para babel y citas

\usepackage{import}
% mode=image: Usa el PDF pre-compilado, nunca intenta recompilar.
% Debes compilar cada figura manualmente antes de compilar el libro.
\usepackage[mode=image, subpreambles=false]{standalone}

% --- 2. MATEMÁTICAS Y CIENCIAS ---
\usepackage{amsmath, amssymb, amsfonts, mathtools}
\usepackage{enumitem}

% LaTeX trae \DeclareMathSizes{5}{5}{5}{5}, así que a 5 pt (\tiny) los subíndices
% salen del mismo tamaño que la base: en las etiquetas de figuras el M de
% $\omega_M$ queda desproporcionado. Se restablece la proporción habitual.
\DeclareMathSizes{5}{5}{3.7}{3}

% --- 3. DISEÑO Y GEOMETRÍA --- 
% Unificamos los márgenes para todo el documento
\usepackage[left=2.5cm,right=2.5cm,top=3cm,bottom=3cm]{geometry}
\makeatletter
\ifstandalone % Evita que geometry dañe el recorte de las figuras bajándolas 3cm
  \@ifclasswith{standalone}{crop=false}{
    % Es un capítulo/sección: Dejamos que geometry actúe.
  }{
    % Es una figura: Neutralizamos geometry para que no las recorte mal.
    \geometry{pass}
  }
\fi
\makeatother
\usepackage{parskip} % Espaciado entre párrafos sin sangría
\usepackage{float}   % Para usar [H] en figura

% Ajustes de flotantes para evitar espacios verticales exagerados
\renewcommand{\topfraction}{0.95}      % Permite que las figuras ocupen hasta el 95% superior
\renewcommand{\bottomfraction}{0.95}   % Permite que las figuras ocupen hasta el 95% inferior
\renewcommand{\textfraction}{0.05}     % Permite hojas con tan solo un 5% de texto
\renewcommand{\floatpagefraction}{0.8} % Requiere que una página de solo figuras esté 80% llena
\setcounter{topnumber}{4}
\setcounter{bottomnumber}{4}
\setcounter{totalnumber}{10}

% --- 4. GRÁFICOS Y TABLAS ---
\usepackage{graphicx}
\usepackage{booktabs} % Tablas profesionales
\usepackage{caption}  % Control de captions y \captionof
\usepackage{tikz}
\usetikzlibrary{arrows.meta, calc, angles, quotes, babel, backgrounds}
\usepackage{pgfplots}
\usepgfplotslibrary{groupplots}
\usepgfplotslibrary{external}
\usepgfplotslibrary{patchplots}
\pgfplotsset{compat=1.18}
\usepackage[american, straightvoltages]{circuitikz}

% --- 5. COLORES Y CAJAS ---

\usepackage[table]{xcolor}
\usepackage{tcolorbox}

% Definición de colores corporativos/estilo
\definecolor{utnblue}{RGB}{0, 51, 102}
\definecolor{codegreen}{rgb}{0,0.6,0}
\definecolor{codegray}{rgb}{0.5,0.5,0.5}
\definecolor{codepurple}{rgb}{0.58,0,0.82}
\definecolor{backcolour}{rgb}{0.96,0.96,0.96}

% --- 6. ENTORNOS PERSONALIZADOS (CAJAS) ---

% Caja "Resultado" (Azul Corporativo) — Definiciones, fórmulas clave, resultados importantes.
% Uso: \begin{resultado}[Fórmula de Euler] ... \end{resultado}
\newtcolorbox{resultado}[1][]{
    colback=utnblue!5!white,
    colframe=utnblue,
    title=\textbf{#1},
    fonttitle=\bfseries,
    sharp corners,
    boxrule=0.5mm
}

% Caja "Nota" (Gris) — Advertencias, aplicaciones, contexto secundario.
% Uso: \begin{nota}[Cuidado con la calculadora] ... \end{nota}
% Uso: \begin{nota} ... \end{nota}  → título por defecto "Nota"
\newtcolorbox{nota}[1][Nota]{
    colback=codegray!5!white,
    colframe=codegray,
    title=\textbf{#1},
    fonttitle=\bfseries,
    sharp corners,
    boxrule=0.5mm
}

% Caja inline para resaltar un valor y enlazarlo con su otra aparición en una tabla
% o en una cuenta: mismo color, mismo valor.
% Uso: \marcavalor{$4$}  o  \marcavalor[codegreen]{$4$}
\newtcbox{\marcavalor}[1][utnblue]{
    on line,
    boxsep=1pt, left=2pt, right=2pt, top=1pt, bottom=1pt,
    colback=#1!10!white,
    colframe=#1,
    boxrule=0.4pt,
    arc=2pt
}

% --- 7. CÓDIGO FUENTE (LISTINGS) ---
\usepackage{listings}

\lstdefinestyle{miestilo}{
    backgroundcolor=\color{backcolour},   
    commentstyle=\color{codegreen},
    keywordstyle=\color{magenta},
    numberstyle=\tiny\color{codegray},
    stringstyle=\color{codepurple},
    basicstyle=\ttfamily\footnotesize,
    breakatwhitespace=false,         
    breaklines=true,                 
    captionpos=b,                    
    keepspaces=true,                 
    numbers=left,                    
    numbersep=5pt,                  
    showspaces=false,                
    showstringspaces=false,
    showtabs=false,                  
    tabsize=2,
    frame=single,                    
    rulecolor=\color{codegray}
}

\lstset{style=miestilo}
% Soporte para tildes y ñ en bloques de código
\lstset{literate={ñ}{{\~n}}1 {á}{{\'a}}1 {é}{{\'e}}1 {í}{{\'i}}1 {ó}{{\'o}}1 {ú}{{\'u}}1}

% Operadores matemáticos personalizados

% Vector en sentido discreto (lista finita de coordenadas): negrita + flecha.
% Uso: \vect{v}, \vect{e}_k, \vect{0}.
\newcommand{\vect}[1]{\vec{\mathbf{#1}}}

% Fecha de versión y comando para capítulos standalone
\providecommand{\verdate}{\the\year.\ifnum\month<10 0\fi\the\month.\ifnum\day<10 0\fi\the\day}
\newcommand{\chapterversion}{%
    \vspace{-1.5cm}\begin{flushright}\small\textit{\color{codegray}Versión~\verdate}\end{flushright}\vspace{0.5cm}%
}

% --- 8. HIPERVÍNCULOS (DEBE SER EL ÚLTIMO PAQUETE) ---
\usepackage[colorlinks=true, linkcolor=utnblue, urlcolor=utnblue, citecolor=utnblue]{hyperref}

% --- 9. REFERENCIAS CRUZADAS ENTRE CAPÍTULOS STANDALONE ---
% Cuando se compila un capítulo standalone, xr-hyper lee los .aux de los
% otros capítulos (si existen) para resolver \ref{} a labels externos.
% En la compilación del libro completo este bloque no se activa.
\ifstandalone
  \usepackage{xr-hyper}
  \makeatletter
  \newcommand{\xrchap}[1]{%
    \edef\xrc@self{\detokenize{Capitulo#1}}%
    \edef\xrc@job{\jobname}%
    \ifx\xrc@self\xrc@job\else
      \IfFileExists{../#1capitulo/Capitulo#1.aux}%
        {\externaldocument{../#1capitulo/Capitulo#1}}{}%
    \fi
  }
  \makeatother
  \xrchap{01}\xrchap{02}\xrchap{03}\xrchap{04}
  \xrchap{05}\xrchap{06}\xrchap{07}\xrchap{08}
  \xrchap{09}\xrchap{10}\xrchap{11}\xrchap{12}
  \xrchap{13}
\fi
% \PassOptionsToPackage{gray}{xcolor}

% fig_alias_sinusoide
% Ejemplo concreto de §9.1 con valores de frecuencia angular:
%   coseno de 9 ciclos/s  -> w0 = 18 pi rad/s
%   muestreo de 8 por seg -> ws = 16 pi rad/s, ws/2 = 8 pi rad/s
%   alias                 -> wa = |18pi - 16pi| = 2 pi rad/s (1 ciclo/s)
% Unidad del eje: 1 unidad de dibujo = 2 pi rad/s.
% (a) lineas del coseno en +-18pi, las dos fuera de la banda base.
% (b) replicas a paso ws=16pi; las que caen en banda base (+-2pi) en rojo.
%     Replicas tenues en +-18pi (originales) y +-14pi (de la familia opuesta).
% Banda base sombreada en gris; eje de ordenadas y origen marcados.

\begin{document}
\begin{tikzpicture}[>={Latex[length=4pt]}, font=\scriptsize, xscale=0.45,
    lin/.style={utnblue, thick, -{Latex[length=4pt]}},
    linf/.style={utnblue!40, thick, -{Latex[length=3.5pt]}},
    alias/.style={red!80!black, very thick, -{Latex[length=4pt]}},
    nyq/.style={gray!70, thin, dashed},
    ax/.style={->, thin, black},
    tick/.style={thin},
]
    \def\H{0.9}

    % ===================== (a) coseno antes de muestrear =====================
    \begin{scope}[shift={(0,0)}]
        \node at (-11.6,0.7) {(a)};
        % banda base
        \fill[gray!10] (-4,0) rectangle (4,\H+0.10);
        \draw[ax] (-11,0) -- (11.4,0) node[right] {$\omega$};
        % eje de ordenadas y origen
        \draw[ax] (0,0) -- (0,\H+0.30);
        \node[anchor=south east, font=\tiny, inner sep=1.5pt] at (-0.15,0.04) {$0$};
        % bordes de la banda base
        \draw[nyq] (4,0) -- (4,\H+0.10);
        \draw[nyq] (-4,0) -- (-4,\H+0.10);
        \node[gray!70, anchor=south, font=\tiny] at (4,\H+0.10) {$\tfrac{\omega_s}{2}$};
        \node[gray!70, anchor=south, font=\tiny] at (-4,\H+0.10) {$-\tfrac{\omega_s}{2}$};
        \node[gray!65, font=\tiny] at (2.1,0.45) {banda base};
        % lineas del coseno
        \draw[lin] (9,0) -- (9,\H);
        \draw[lin] (-9,0) -- (-9,\H);
        \node[utnblue, anchor=south, font=\tiny] at (9,\H) {$\omega_0$};
        \node[utnblue, anchor=south, font=\tiny] at (-9,\H) {$-\omega_0$};
        % marcas del eje
        \foreach \x/\lab in {4/{8\pi}, 9/{18\pi}, -4/{-8\pi}, -9/{-18\pi} } {
            \draw[tick] (\x,0.05) -- (\x,-0.05) node[below, font=\tiny] {$\lab$};
        }
    \end{scope}

    % ===================== (b) despues de muestrear =====================
    \begin{scope}[shift={(0,-3.0)}]
        \node at (-11.6,0.7) {(b)};
        % banda base
        \fill[gray!10] (-4,0) rectangle (4,\H+0.10);
        \draw[ax] (-11,0) -- (11.4,0) node[right] {$\omega$};
        % eje de ordenadas y origen
        \draw[ax] (0,0) -- (0,\H+0.30);
        \node[anchor=south east, font=\tiny, inner sep=1.5pt] at (-0.15,0.04) {$0$};
        % bordes de la banda base
        \draw[nyq] (4,0) -- (4,\H+0.10);
        \draw[nyq] (-4,0) -- (-4,\H+0.10);
        \node[gray!70, anchor=south, font=\tiny] at (4,\H+0.10) {$\tfrac{\omega_s}{2}$};
        \node[gray!70, anchor=south, font=\tiny] at (-4,\H+0.10) {$-\tfrac{\omega_s}{2}$};
        % replicas fuera de la banda base (tenues)
        \draw[linf] (9,0) -- (9,\H);
        \draw[linf] (-9,0) -- (-9,\H);
        \draw[linf] (7,0) -- (7,\H);
        \draw[linf] (-7,0) -- (-7,\H);
        % alias dentro de la banda base
        \draw[alias] (1,0) -- (1,\H);
        \draw[alias] (-1,0) -- (-1,\H);
        \node[red!80!black, anchor=south west, font=\tiny, inner sep=1pt]
            at (1,\H) {$\omega_a$};
        % marcas del eje
        \foreach \x/\lab in {1/{2\pi}, 4/{8\pi}, 9/{18\pi},
                             -1/{-2\pi}, -4/{-8\pi}, -9/{-18\pi} } {
            \draw[tick] (\x,0.05) -- (\x,-0.05) node[below, font=\tiny] {$\lab$};
        }
        \foreach \x/\lab in {7/{14\pi}, -7/{-14\pi} } {
            \draw[tick, gray!60] (\x,0.05) -- (\x,-0.05)
                node[below, font=\tiny, gray!70] {$\lab$};
        }
        % frecuencia de muestreo sobre el eje
        \draw[tick, gray!70] (8,0.05) -- (8,-0.05);
        \draw[tick, gray!70] (-8,0.05) -- (-8,-0.05);
        \node[gray!70, anchor=north, font=\tiny] at (8,-0.42) {$\omega_s\!=\!16\pi$};
        \draw[gray!70, thin, -{Latex[length=3pt]}] (8,-0.42) -- (8,-0.13);
        \node[gray!70, anchor=north, font=\tiny] at (-8,-0.42) {$-\omega_s$};
        \draw[gray!70, thin, -{Latex[length=3pt]}] (-8,-0.42) -- (-8,-0.13);
        % traslado de +-18pi hasta +-2pi
        \draw[red!80!black, thin, -{Latex[length=3pt]}]
            (8.7,\H+0.42) to[bend right=14] (1.3,\H+0.42);
        \node[red!80!black, anchor=south, font=\tiny] at (5,\H+0.52) {$-\,\omega_s$};
        \draw[red!80!black, thin, -{Latex[length=3pt]}]
            (-8.7,\H+0.42) to[bend left=14] (-1.3,\H+0.42);
        \node[red!80!black, anchor=south, font=\tiny] at (-5,\H+0.52) {$+\,\omega_s$};
    \end{scope}
\end{tikzpicture}
\end{document}
